\subsection{Finite formula families and infinitesimal approximation}
We describe the local topology of a finite definable family of
homeomorphisms and the separate limitation of finite PL approximation.

Let $R$ be a real closed field with a fixed o-minimal expansion,
and let $K\subseteq R^n$ be closed, bounded, and definable.
Fix finitely many formulas for graphs of maps $K\to K$, with
finite parameter tuples.  Retain the parameters whose graphs are
homeomorphisms satisfying a prescribed finite list of definable
marked-set or boundary conditions.  Identify parameters that define
the same whole map, also across different formulas, and write
$\mathcal H$ for the resulting definable quotient.  No closure under
composition or inversion within the chosen formats is assumed.
For nonempty $K$, give it the metric
\begin{equation}\label{eq:finite-format-uniform-metric}
 d(h,g)=\max_{x\in K}\|h(x)-g(x)\|
          +\max_{x\in K}\|h^{-1}(x)-g^{-1}(x)\|.
\end{equation}
All topologies and inequalities here are over $R$, including when
$R$ is non-Archimedean.  The following application uses Walsberg's
affineness theorem for definable metric spaces
\cite[Theorem~7.1, p.~16]{Walsberg2015}.

\begin{theorem}\label{thm:finite-format-homeomorphism-contractions}
The metric space $(\mathcal H,d)$ is definably homeomorphic to a
definable subset of some $R^N$ with its Euclidean topology.
For every $h_0\in\mathcal H$ and every $\epsilon>0$, there are a
definable $d$-open neighborhood $U$ of $h_0$, a $\delta>0$, and a
definable continuous contraction $C:U\times I\to U$ satisfying
\[
 B_d(h_0,\delta)\subseteq U\subseteq B_d(h_0,\epsilon),\qquad
 C(h,0)=h,\quad C(h,1)=h_0,\quad C(h_0,s)=h_0.
\]
Both $(h,s,x)\mapsto C(h,s)(x)$ and
$(h,s,x)\mapsto C(h,s)^{-1}(x)$ are jointly definable and continuous.
Thus the contraction consists of actual homeomorphisms in the same
finite family of formats satisfying the prescribed marks.  If $K$
is empty, the quotient is empty or a singleton and the assertion is
interpreted accordingly.
\end{theorem}

\begin{proof}
Admissibility of a parameter tuple is definable: graph uniqueness,
bijectivity, continuity, inverse continuity, and the finite marked
requirements are first-order conditions.  A finite label combines
the different formats into a single definable parameter space.
Equality of their whole graphs is a definable equivalence relation.
Elimination of imaginaries for o-minimal expansions of fields
realizes its quotient as a definable set; this is also the quotient
convention in \cite[p.~1]{Walsberg2015}.  Evaluation and inverse
evaluation are definable on that quotient.  Definable compactness
of $K$ gives the two maxima in
\eqref{eq:finite-format-uniform-metric}; their sum is a definable
metric.

We verify the hypothesis of Walsberg's theorem over the given
field: every definable $d$-discrete subset is finite.  Suppose
instead that $D\subseteq\mathcal H$ is an infinite definable
$d$-discrete subset.  Apply definable Hardt triviality to the
family of its compact graphs and refine the parameter partition
into cells \cite[Theorem~5.22, pp.~62--63]{Coste1999}.  One cell
$A\subseteq D$ has positive dimension, and over it there is a
definable graph trivialization
\[
 \theta:A\times G\longrightarrow
 \{(a,x,h_a(x)):a\in A,\ x\in K\},
\]
where $G$ is one closed bounded graph fibre.  Projection after
$\theta$ to $(a,x)$ is a continuous definable bijection
$A\times G\to A\times K$.

This bijection is locally a homeomorphism on the parameter cell.
Indeed, a cell is locally closed, and each of its points has a
closed bounded definable neighborhood $L$ within the cell
\cite[Proposition~3.7 and Exercise~3.8, pp.~27--28]{Coste1999}.
On $L\times G$ the projection is a continuous definable bijection
between closed bounded definable sets, so its inverse is continuous.
One can check this without ordinary compactness: the image of
every closed definable subset is closed by definable compactness,
which proves inverse continuity on definable basic neighborhoods.
Projection to $(a,y)$ gives the same conclusion for the inverse
graphs.  Consequently both $(a,x)\mapsto h_a(x)$ and
$(a,y)\mapsto h_a^{-1}(y)$ are jointly continuous on $A\times K$.

On each $L\times K$ these maps are uniformly continuous, by
definable compactness.  For example, for a continuous definable
map on a closed bounded set, the pairs of points whose images
differ by at least $\epsilon$ form a closed bounded set disjoint
from the diagonal.  If nonempty, the distance of its points from
the diagonal has a positive minimum.  This supplies a uniform
modulus at $\epsilon$; the empty case is immediate.  Applied to
the two evaluation maps, it shows that
$a\mapsto h_a$ is continuous from the Euclidean cell $A$ to
$(\mathcal H,d)$.  A sufficiently small $d$-ball about any $a\in A$
meets $D$ only in $a$.  Its inverse image would make $\{a\}$ open
in the positive-dimensional cell, which is impossible.  This
contradiction proves definable separability in the exact sense
used by Walsberg.  Theorem~7.1 therefore gives a definable
homeomorphism $\varphi:(\mathcal H,d)\to X\subseteq R^N$.
No ordinary separability or Archimedean hypothesis has been used.

We spell out the strong local contraction, including the possibility
that $X$ is not locally closed.  Replace the ambient Euclidean
space by a bounded definably homeomorphic copy if necessary.
Put $x_0=\varphi(h_0)$ and
$V=\varphi(B_d(h_0,\epsilon))$.  Triangulate the closed bounded
set $\overline X$ compatibly with $X$, $V$, and $\{x_0\}$
\cite[Theorem~4.4, pp.~42--44]{Coste1999}.  The inverse image of
$x_0$ is a vertex $v$.  In the triangulation, intersect the open
star of $v$ with the inverse image of $X$.  This is an open
neighborhood of $v$ in that inverse image.  It is contained in the
inverse image of $V$: otherwise a compatible open simplex lying
in $X\setminus V$ would have $v$ as a face, contrary to the
relative openness of $V$ at $x_0$.

Straight contraction to $v$ preserves this neighborhood.  For a
point in an open simplex containing $v$ as a face, its segment to
$v$ remains in that open simplex until the endpoint.  Thus no
missing faces of $X$ are introduced.  Transport the contraction
back through the triangulation and $\varphi$.  This gives $C$ and
$U$ with the stated endpoint identities and with $U$ invariant
throughout the contraction.  As $U$ is $d$-open at $h_0$, it contains
a ball of some positive radius $\delta$.

Finally, evaluation is continuous for the metric topology directly:
at $(g,y)$,
\[
 \|h(x)-g(y)\|
 \leq\max_{z\in K}\|h(z)-g(z)\|+\|g(x)-g(y)\|.
\]
The first term is controlled by $d(h,g)$ and the second by
continuity of $g$.  Applying the same argument to inverses proves
joint inverse evaluation.  Composition with $C$ proves both
family assertions.  No continuous choice of coefficient tuples
is needed.
\end{proof}

\begin{proposition}\label{prop:infinitesimal-pl-approximation-obstruction}
If $p:[0,1]\to R$ is continuous and affine on at most $m$ intervals,
where $m\geq1$ is an ordinary finite integer, then
\begin{equation}\label{eq:finite-piece-quadratic-lower-bound}
 \max_{x\in[0,1]}|x^2-p(x)|\geq\frac{1}{8m^2}.
\end{equation}
The coefficients and breakpoints may be arbitrary elements of $R$.
Hence no finite PL map approximates $x\mapsto x^2$ to a positive
infinitesimal uniform tolerance over a non-Archimedean real closed
field.  Nevertheless $x\mapsto x^2$ is semialgebraically isotopic
to the identity relative to the endpoints.
\end{proposition}

\begin{proof}
One affine interval $[a,b]$ has length at least $1/m$.  Put
$c=(a+b)/2$ and $e(x)=x^2-p(x)$.  Affineness gives
\[
 e(a)+e(b)-2e(c)=\frac{(b-a)^2}{2}.
\]
If $E$ is the uniform error, the left-hand side has absolute value
at most $4E$.  Thus $E\geq(b-a)^2/8\geq1/(8m^2)$.
The argument is ordered-field arithmetic, so allowing nonstandard
coefficients or breakpoints does not change it.  A positive
infinitesimal is smaller than each of these positive rational bounds
as $m$ ranges over ordinary finite integers.

The semialgebraic family
\[
 H_s(x)=(1-s)x^2+sx,\qquad 0\leq s\leq1,
\]
fixes both endpoints.  If $x<y$ in $[0,1]$, then
\[
 H_s(y)-H_s(x)=(y-x)\bigl((1-s)(x+y)+s\bigr)>0.
\]
Each slice is therefore an increasing semialgebraic homeomorphism
of the interval.  The total map $(x,s)\mapsto(H_s(x),s)$ is a
continuous definable bijection of the closed bounded square, so
definable compactness makes its inverse continuous.  It is the
required genuine semialgebraic isotopy.
\end{proof}

The contraction neighborhood depends on the finite formula family.
No contraction on the union of all formats is asserted.  The
infinitesimal approximation obstruction concerns finite PL
approximation and does not contradict the displayed isotopy or
establish a failure of a higher concordance comparison.
